\documentclass[10pt,a4paper]{amsart}
\usepackage[margin=19mm]{geometry}
\usepackage{amsmath,amssymb,mathtools}
\usepackage[scr=rsfs,cal=euler]{mathalfa}
\usepackage{xfrac}
\usepackage[unicode,hidelinks,bookmarks=false]{hyperref}
\newcommand{\ie}{i.e.\ }
\newcommand{\cf}{cf.\ }
\newcommand{\resp}{respectively\ }
\newcommand{\Subsection}{\subsection}
\providecommand{\nobreakdash}{\nobreak\hbox{-}}
\setlength{\emergencystretch}{2em}
\multlinegap0pt
\usepackage{bm}
\usepackage{bbm}
\usepackage{dsfont}
\usepackage{xspace}
\usepackage{cancel}
\usepackage{mathtools}
\usepackage{amsthm}
\makeatletter
\@ifundefined{theorem}{\newtheorem{theorem}{Theorem}}{}
\makeatother
\newcommand{\lx}{\mathfrak{L}}
\title[Lexical tableaux and quasisymmetric functions]{Lexical tableaux and quasisymmetric functions}
\author{John M. Campbell, Spencer Daugherty}
\date{Source: arXiv:2511.00713v1 (CC BY 4.0)}
\begin{document}
\maketitle

\noindent\emph{Source and licence.} Lexical tableaux and quasisymmetric functions. Version arXiv:2511.00713v1 (CC BY 4.0), distributed under the Creative Commons Attribution 4.0 International licence, as recorded in the arXiv licence metadata for this exact version and shown on the abstract page https://arxiv.org/abs/2511.00713v1. Full rights notice: licenses/arxiv-2511.00713-CC-BY-4.0.txt.\par\smallskip

\noindent\textit{Editorial scope.} Counted unit: the Theorem whose source label is \texttt{theorem:lexical\_coeff} in the article (source0.tex lines 553--555, source label \texttt{theorem:lexical\_coeff}), delivered together with the article's own complete original proof of it (source0.tex lines 557--573, one \texttt{proof} environment of 17 lines). No mathematics is written, repaired or re-derived: statement and proof are the article's own text line for line. Required context retained uncounted: none. Every definition, display and label of those lines is retained unchanged. Omitted from this package and not redistributed: 1--552, 556--556, 574--1113. Nothing inside a retained excerpt is omitted or altered: every retained body is delivered exactly as the article's lines are written, so no omission receipt of any kind accompanies this package. Citations: the retained excerpts carry no citation command at all, so no bibliography entry of the article is transcribed and no reference is redistributed. Reference adaptation: none was needed and none was made; every reference inside the delivered unit resolves inside it, so no unresolved reference or citation is produced. Printed slips kept literal: no printed word, symbol, hypothesis or number of the counted statement or of its proof was altered. Layout and class: the article's own class is \texttt{amsart} with packages including geometry, natbib, shuffle, mathrsfs, bbm, ifthen, tikz, mathdots, array, multirow, cjhebrew, hyperref, float, caption; the wrapper is the lane's pinned \texttt{amsart} editorial wrapper at 10pt on A4, which restates the theorem-like environments the retained text needs and adds the article's own macro definitions that the retained text needs (\texttt{\textbackslash lx}), with extra wrapper packages (bm, bbm, dsfont, xspace, cancel, mathtools). The delivered numbering is the unit's own. Assets: no retained span contains a figure environment, a graphics-inclusion command, a PDF-page inclusion, a table or a TikZ picture; no image file is copied into this package, the package declares no asset dependency and no image file is redistributed with it. Licence: version arXiv:2511.00713v1 of this article is distributed under the Creative Commons Attribution 4.0 International licence, as recorded in the arXiv licence metadata for this exact version and shown on the abstract page https://arxiv.org/abs/2511.00713v1. A full rights notice is delivered at licenses/arxiv-2511.00713-CC-BY-4.0.txt. Characters: every retained excerpt is pure ASCII, so the delivered engine, XeLaTeX with its default Latin Modern text font, reports no missing character.
\begin{theorem}\label{theorem:lexical_coeff}
 Given a composition $\alpha \vDash n$, we have $K^{\lx}_{\alpha, \alpha}=1$ and, if $\alpha \leq_{\ell} \beta$, then $K^{\lx}_{\alpha, \beta}=0$.
\end{theorem}

\begin{proof}
 Consider an empty diagram of shape $\alpha = (\alpha_1, \ldots, \alpha_k)$ that we want to fill as a lexical tableau of type $\alpha$. If 
 there is a $1$ anywhere in a given row, then the first entry of that row must be a $1$. Since the first column is strictly increasing, there
 must be a $1$ in the first row, and there must not be any $1$s in the following rows. So all $\alpha_1$ of the $1$s in the tableaux must 
 go in the top row, and they fill it completely. Next, we need to fill $\alpha_2$ cells with $2$. Here, $2$ will be the lowest entry in any of 
 the rows (other than the first), so any row with a $2$ must have a $2$ as its first entry. Since the first column is strictly increasing, 
 the second row must have a $2$ as its first entry, and no other row may contain any $2$'s. Therefore, we fill the second row entirely 
 with $2$'s, which uses all $\alpha_2$ that we seek to use. Continuing this pattern, we see that the only way to construct a lexical tableau 
 of shape $\alpha$ and type $\alpha$ is to fill each of the $\alpha_i$ cells of row $i$ with $i$'s, and thus  $K_{\alpha,\alpha} = 1$. 
 
 Next, consider some $\alpha \leq_{\ell} \beta$, meaning there exists some $j$ such that $\alpha_1 = \beta_1, \alpha_2 = \beta_2, \ldots, 
 \alpha_{j-1}=\beta_{j-1},$ and $\alpha_j<\beta_j$. By way of contradiction, suppose that there exists a lexical tableau $T$ of shape $\alpha$ and type $\beta$. Using a similar argument, relative to our preceding argument, the first $j-1$ rows of $T$ are necessarily filled entirely with the integer that matches their row index. That is, if $i < j$ then row 
 $i$ is has exactly $\alpha_i = \beta_i$ each filled with an $i$. Next, we will fill in the $\beta_j$ instances of $j$. Given the current filling, 
 any row with a $j$ must have a $j$ as its first entry. Since the first column is strictly increasing, row $j$ must contain all of these $j$'s. 
 However, there are only $\alpha_j$ cells in row $j$ and we need to place $\beta_j > \alpha_j$ instances of $j$. Thus, we cannot create a 
 lexical tableau of shape $\alpha$ and type $\beta$, meaning $K^{\lx}_{\alpha, \beta}=0$ if $\alpha \leq_{\ell} \beta$. 
\end{proof}

\end{document}
